\chapter{Historia de la responsabilidad extracontractual del Estado}
\section{Evolución de la responsabilidad del Estado}
\begin{itemize}
\item \textbf{Edad Media:} \emph{The king can do no wrong}; inmunidad o irresponsabilidad del gobernante, especialmente durante los regímenes absolutistas, bajo la concepción de que el soberano no respondía por los daños causados en ejercicio del poder.
\item \textbf{Surgimiento del principio de legalidad y del Estado de derecho:} también las autoridades están sometidas al ordenamiento jurídico. El principio de legalidad tiene dos dimensiones: los particulares pueden hacer todo aquello que no esté prohibido por una norma previa y cierta; las autoridades solo pueden actuar dentro de las competencias que la Constitución y la ley les atribuyen previamente. A partir de allí se fortalece la posibilidad de exigir responsabilidad al Estado cuando su actuación causa daños.
\item \textbf{Responsabilidad extracontractual del Estado basada en normas del derecho civil:} inicialmente se analizaba aplicando reglas del Código Civil.
\item \textbf{Autonomía de la responsabilidad extracontractual del Estado:} progresivamente se construyó un régimen autónomo, con principios, normas y criterios propios del derecho administrativo.
\end{itemize}
\subsection{Antecedentes: la Escuela del Servicio Público}
En Francia adquirió especial importancia la Escuela del Servicio Público, según la cual el derecho administrativo encontraba buena parte de su justificación en que una de las funciones centrales del Estado era la prestación de servicios públicos orientados a satisfacer el interés general. Desde esta perspectiva, la actividad estatal no podía regirse exclusivamente por el derecho privado, pues el Estado actuaba con finalidades y potestades distintas de las de los particulares. Se consolidó así la idea de un régimen jurídico especial y autónomo, propio del derecho administrativo, aplicable tanto a la organización y prestación de los servicios públicos como a la responsabilidad derivada de su funcionamiento.

\subsection{El surgimiento de la jurisdicción contencioso administrativa en Francia}
Después de la Revolución francesa se consolidó una fuerte desconfianza frente a los jueces ordinarios, porque durante el Antiguo Régimen algunos tribunales habían obstaculizado reformas impulsadas por el poder político. Por ello, inicialmente se limitó la posibilidad de que los jueces intervinieran en asuntos administrativos y se fortalecieron órganos con funciones consultivas dentro de la propia Administración. Con el tiempo, Francia desarrolló un modelo de justicia dual: una jurisdicción ordinaria para los asuntos privados y una jurisdicción administrativa especializada para los conflictos relacionados con la actuación del Estado. Este modelo contrasta con sistemas de justicia más unitarios, como el de Estados Unidos.

\section{El Fallo Blanco (Tribunal de Conflictos, 8 de febrero de 1873)}
\begin{cita}«Considerando que la responsabilidad, que podría incumbir al Estado por los daños ocasionados a particulares, como resultado de la actuación de personas empleadas en los servicios públicos, no puede regirse por los principios que se establecen en el Código Civil para las relaciones entre particulares. Que dicha responsabilidad no es ni general ni absoluta; que se rige por normas especiales que varían en función de las necesidades del servicio y de la necesidad de conciliar los derechos del Estado con los derechos privados; Que, por tanto, en virtud de las leyes antedichas, la autoridad administrativa es la única competente para conocer de esta causa;...»\end{cita}
\begin{ejemplo}Los hechos. Agnès Blanco, una niña, fue atropellada y gravemente lesionada por un vagón perteneciente a una fábrica de tabacos administrada por el Estado francés. Su padre demandó buscando que se aplicaran las reglas del Código Civil sobre responsabilidad.\end{ejemplo}
El Tribunal de Conflictos adoptó dos decisiones fundamentales:

\begin{itemize}
\item \textbf{Decisión procesal:} el asunto debía ser conocido por la jurisdicción administrativa, porque el daño estaba relacionado con el funcionamiento de un servicio público estatal.
\item \textbf{Decisión sustancial:} la responsabilidad del Estado por los daños causados por sus servidores no puede regirse automáticamente por las normas del Código Civil aplicables entre particulares, porque están de por medio la prestación de un servicio público, el interés general y la naturaleza especial de la actividad estatal. No es ni general ni absoluta: se rige por normas especiales propias del derecho administrativo, que varían según las necesidades del servicio y la necesidad de conciliar los derechos del Estado con los derechos privados.
\end{itemize}
\textbf{Importancia.} Es considerado una piedra angular del derecho administrativo moderno, porque consolidó la autonomía de la responsabilidad del Estado frente al derecho civil y reforzó la existencia de una jurisdicción especializada para conocer de estos conflictos.

\begin{comentario}Este fallo explica también la tensión entre las necesidades del servicio y los derechos de las víctimas, que está en la base del \emph{favor victimae} (ver Capítulo 7, sección 7.6).\end{comentario}
\section{La responsabilidad del Estado en Colombia}
\subsection{Nacimiento de la jurisdicción contencioso administrativa en Colombia}
\begin{itemize}
\item \textbf{Origen consultivo.} De manera semejante al antecedente francés, el Consejo de Estado tuvo inicialmente funciones consultivas. Fue creado por Simón Bolívar mediante un decreto expedido en Angostura el 30 de octubre de 1817 (no por el Congreso de Angostura, que se reunió en 1819). Durante el siglo XIX su existencia fue intermitente: algunas constituciones lo contemplaron y otras lo suprimieron.
\item \textbf{La Constitución de 1886} permitió establecer la jurisdicción contencioso administrativa: su artículo 164 habilitó al legislador para organizarla y el artículo 141 contempló que el Consejo de Estado ejerciera funciones jurisdiccionales si la ley establecía esa jurisdicción.
\item \textbf{La Ley 130 de 1913} materializó su organización y constituye el primer Código Contencioso Administrativo. En 1914 se restableció el Consejo de Estado como órgano consultivo y tribunal supremo de esta jurisdicción.
\item \textbf{Dualidad de jurisdicciones:} se adoptó la separación entre la jurisdicción ordinaria y la contencioso administrativa, siguiendo la influencia francesa y la necesidad de contar con un juez especializado en asuntos administrativos.
\item \textbf{Competencia inicial y diferencia con Francia:} la jurisdicción contencioso administrativa colombiana comenzó principalmente como juez de legalidad de los actos administrativos, mediante mecanismos antecedentes de las actuales pretensiones de nulidad y de nulidad y restablecimiento del derecho; inicialmente no tenía una competencia general sobre responsabilidad contractual y extracontractual del Estado. En Francia, la jurisdicción administrativa ya había desarrollado competencias sobre responsabilidad estatal (aunque no conocía todos los pleitos del Estado, pues también había asuntos atribuidos a la justicia ordinaria). Por esa distribución, buena parte de las controversias sobre responsabilidad extracontractual estatal en Colombia fue resuelta inicialmente por la jurisdicción ordinaria.
\item \textbf{Ampliación de competencias:} la Ley 167 de 1941 amplió las competencias de la jurisdicción administrativa en asuntos específicos, como las indemnizaciones por trabajos públicos. El traslado general de las controversias sobre contratos administrativos y responsabilidad extracontractual se consolidó con el Decreto 528 de 1964 (no ocurrió de manera general en 1941 ni mediante decretos leyes de ese año).
\end{itemize}
\subsection{Períodos de la responsabilidad del Estado en Colombia}
{\small\begin{xltabular}{\linewidth}{|X|X|}\hline
\rowcolor{navy!12}\textbf{Período} & \textbf{Jurisdicción competente} \\\hline\endhead
1886-1941 & Jurisdicción ordinaria \\\hline
1941-1991 & Jurisdicción de lo contencioso administrativo \\\hline
\end{xltabular}}
\subsection{Fundamentos aplicados por la jurisdicción ordinaria (posiciones de la Corte Suprema de Justicia)}
La Corte Suprema construyó inicialmente la responsabilidad estatal aplicando las reglas del Código Civil. Su evolución puede organizarse en etapas, aunque las distintas tesis coexistieron y no se sustituyeron de forma completamente sucesiva.

\textbf{1. Primer período – Personas jurídicas: remisión a las normas sobre responsabilidad indirecta.} La Corte aplicó al Estado la responsabilidad por el hecho ajeno que utilizaba para las personas jurídicas: el funcionario causaba el daño y la entidad respondía por la culpa en su elección (\emph{in eligendo}) o vigilancia (\emph{in vigilando}). Esta tesis podía desfavorecer a la víctima por dos razones:

\begin{itemize}
\item El término de prescripción de tres años aplicable a la responsabilidad indirecta era más corto.
\item La discusión se centraba en la elección o vigilancia del agente (esta culpa podía presumirse, por lo que no siempre correspondía a la víctima demostrar directamente que la entidad había elegido o vigilado mal).
\end{itemize}
\textbf{2. Segundo período – Personas jurídicas: remisión a las normas sobre responsabilidad directa.} La actuación de los agentes empezó a atribuirse directamente a la persona jurídica, que respondía por una conducta jurídicamente propia.

\begin{itemize}
\item Se aplicaba el término ordinario de prescripción de veinte años vigente en esa época, más amplio y favorable para reclamar.
\item El fundamento dejó de ser la culpa en la elección o vigilancia y pasó a ser la conducta atribuible a la propia entidad.
\end{itemize}
\textbf{3. Tercer período – Tesis organicista.} Distinguía entre dos clases de agentes dentro de la persona jurídica:

\begin{itemize}
\item \textbf{Órganos (agentes órganos) de dirección o representación:} expresaban la voluntad de la entidad y comprometían su responsabilidad directa.
\item \textbf{Agentes subalternos:} ejecutaban las órdenes de los directivos y sus actuaciones generaban responsabilidad indirecta.
\end{itemize}
\begin{ejemplo}El fallo Tinjacá (1962) consolidó la responsabilidad directa y superó esa distinción jerárquica: la actuación funcional del agente se atribuye a la Administración, independientemente de que sea directivo o subalterno.\end{ejemplo}
\textbf{4. Falla del servicio.} La Corte Suprema la desarrolló inicialmente y luego el Consejo de Estado la retomó y profundizó. Con la ampliación de sus competencias, especialmente en 1964, la jurisdicción administrativa consolidó el estudio de la responsabilidad estatal mediante distintos títulos (falla del servicio, daño especial y riesgo excepcional), que tienen antecedentes propios y no surgieron todos después de una transferencia general de competencias en 1941.

\subsection{Fundamentos aplicados por la jurisdicción de lo contencioso administrativo antes de 1991}
Antes de la Constitución de 1991 se aplicaban tres fundamentos: la falla del servicio, el riesgo excepcional y el daño especial.

\paragraph{a) Falla del servicio}
\begin{cita}«ARTÍCULO 16. Las autoridades de la República están instituidas para proteger a todas las personas residentes en Colombia, en sus vidas, honra y bienes, y para asegurar el cumplimiento de los deberes sociales del Estado y de los particulares.» (Constitución de 1886)\end{cita}
\textbf{Fundamento.} El artículo 16 de la Constitución de 1886 establecía el deber de las autoridades de proteger la vida, honra y bienes de las personas. Cuando el Estado incumple ese deber de protección, puede configurarse una falla del servicio.

\textbf{Modalidades.} La falla del servicio se presenta cuando el Estado:

\begin{itemize}
\item \textbf{No presta el servicio:} omisión del servicio.
\item \textbf{Lo presta tardíamente:} actúa después del momento en que debía hacerlo.
\item \textbf{Lo presta defectuosamente:} actúa, pero de manera inadecuada o incumpliendo sus deberes.
\end{itemize}
\textbf{Naturaleza: régimen subjetivo.} Exige demostrar que la Administración incumplió un deber (no prestó el servicio, lo prestó tardíamente o lo prestó defectuosamente). Se examina la conducta de la entidad, sin que sea necesario identificar a un funcionario culpable. Cumple una función equivalente a la culpa en materia civil: se compara cómo actuó la entidad estatal y cómo debía actuar según sus obligaciones y las circunstancias del caso; si su actuación estuvo por debajo del estándar exigible, se configura la falla. La sola ocurrencia del daño no basta: que alguien resulte lesionado o muera no demuestra automáticamente una falla del servicio; debe acreditarse el incumplimiento estatal y su relación con el daño.

\textbf{Características} (relatividad de la falla, falla anónima y falla presunta):

\begin{itemize}
\item \textbf{Relatividad.} La obligación estatal de protección es de medio, no absoluta: el Estado no puede estar omnipresente. Se exige una actuación diligente según las circunstancias, no garantizar que nunca ocurra un daño.
\end{itemize}
\begin{ejemplo}El barrio peligroso. El Estado no puede estar presente en todos los barrios peligrosos del país; por eso, si a alguien le roban en uno de ellos, no existe automáticamente responsabilidad por falla del servicio. Pero si los policías del CAI conocen lo que está ocurriendo y, pudiendo y debiendo intervenir, no hacen nada, puede existir falla del servicio por falta de protección.\end{ejemplo}
\begin{itemize}
\item \textbf{Carácter anónimo.} La demanda se dirige contra la entidad estatal y no es necesario identificar el nombre del agente que causó el daño; basta demostrar la falla atribuible al funcionamiento de la Administración y su relación con el daño.
\item \textbf{Falla presunta y carga dinámica de la prueba.} La regla general es que el demandante prueba los hechos que sustentan su reclamación y el demandado los que fundamentan sus excepciones (regla que estaba en el art. 177 del Código de Procedimiento Civil y hoy está en el art. 167 del Código General del Proceso). En la \textbf{falla presunta}, la falla se presume y corresponde a la entidad desvirtuarla demostrando que actuó adecuadamente; esto no significa que se presuman automáticamente el daño y el nexo causal. La \textbf{carga dinámica de la prueba} (art. 167 del CGP) permite al juez asignar la prueba de un hecho a la parte que esté en mejores condiciones para acreditarlo; lo decide mediante providencia motivada, antes de dictar sentencia, y debe permitir aportar y controvertir la prueba. No se invierte automáticamente toda la carga del proceso, sino que se distribuye respecto de hechos determinados. La carga dinámica y la falla presunta son figuras distintas.
\end{itemize}
\begin{ejemplo}Responsabilidad médica estatal: el Consejo de Estado abandonó la aplicación general de la falla presunta y regresó a la falla probada. Una sentencia clave es la del 31 de agosto de 2006, expediente 15772.\end{ejemplo}
\begin{comentario}El fundamento de la falla del servicio fue el artículo 16 de la Constitución de 1886, que hoy corresponde, en términos bastante similares, al artículo 2.º de la Constitución de 1991 (ver Capítulo 6).\end{comentario}
\paragraph{b) Riesgo excepcional}
\begin{cita}«Tiene ocurrencia cuando el Estado, en desarrollo de una obra de servicio público utiliza recursos o medios que colocan a los particulares o a sus bienes en situación de quedar expuestos a "un riesgo de naturaleza excepcional" (De Laubadère) el cual dada su gravedad, excede las cargas que normalmente deben soportar los mismos particulares como contrapartida de las ventajas que resultan de la existencia de ese servicio público. Si el riesgo llega a realizarse y ocasiona un daño, sin culpa de la víctima, hay lugar a responsabilidad de la Administración, así no haya habido falta o falla del servicio.»\end{cita}
\textbf{Naturaleza: régimen objetivo.} El fundamento de la responsabilidad es la materialización del riesgo, no la culpa de la Administración. Por eso, demostrar que el Estado actuó diligentemente no basta para exonerarlo. Aplica solo cuando el daño proviene de un riesgo.

\textbf{Elementos que deben acreditarse:}

\begin{itemize}
\item Un riesgo excepcional atribuible a la actividad estatal.
\item Un daño antijurídico.
\item La relación entre la materialización del riesgo y el daño.
\end{itemize}
\begin{ejemplo}Caso fundador: la muerte de los novillos (Consejo de Estado, 2 de febrero de 1984, Exp. 2744). Un cable de energía eléctrica cayó sobre un potrero de la finca La Camelia, en Quimbaya; por el cortocircuito murieron seis novillos y otro quedó herido. El Consejo de Estado explicó que la responsabilidad podía surgir del riesgo excepcional creado por la conducción de energía eléctrica, sin necesidad de demostrar una falla del servicio.\end{ejemplo}
\textbf{Casos típicos:} lesiones por disparos accidentales de armas de dotación oficial, electrocuciones causadas por redes eléctricas a cargo de una entidad estatal y vehículos oficiales, siempre que el daño sea atribuible a esa actividad peligrosa. La responsabilidad no es automática: puede existir una causa extraña que exonere al Estado.

\paragraph{c) Daño especial}
\textbf{Concepto.} Se aplica cuando el Estado, obrando lícitamente, causa un daño que rompe el equilibrio de las cargas públicas.

\textbf{Naturaleza: régimen objetivo,} porque no exige demostrar culpa ni falla del servicio: la responsabilidad surge de imponer a la víctima una carga especial y anormal que no tiene el deber jurídico de soportar.

\textbf{Requisitos:}

\begin{itemize}
\item \textbf{Actividad lícita del Estado:} la Administración actúa conforme al ordenamiento jurídico.
\item \textbf{Ruptura de la igualdad ante las cargas públicas:} por vivir en sociedad soportamos cargas e incomodidades ordinarias, como los impuestos o el tráfico, que no generan automáticamente responsabilidad estatal. Pero cuando el Estado, aun actuando lícitamente, impone a una persona o a unas pocas una carga muy superior a la que normalmente soportan los demás, puede generar responsabilidad extracontractual.
\item \textbf{Daño especial y anormal:} la afectación supera las cargas ordinarias de la vida en sociedad y la víctima no está jurídicamente obligada a soportarla.
\item \textbf{Nexo causal} entre la actividad lícita del Estado y el daño.
\end{itemize}
\begin{cita}«De una parte, dice, ningún ciudadano debe sufrir otras cargas que las impuestas a los demás en interés social. De otra, las necesidades de la vida colectiva exigen que cada uno soporte sin indemnización los daños resultantes del ejercicio legal y regular de la potencia pública, a menos que el legislador no disponga otra cosa. Y el segundo así: El daño producido por la Administración no da lugar a reparación sino en el caso en que es anormal por su importancia y por su carácter excepcional. La Administración tiene el derecho de imponer este sacrificio especial como gestora suprema del interés público pero mediante indemnización, al efecto de restablecer la igualdad de las cargas.»\end{cita}
\begin{ejemplo}Caso El Siglo, «el caso de los censores» (Consejo de Estado, 29 de julio de 1947): turbación del orden público por la captura del Presidente en un alzamiento militar en Pasto. Se omitió designar al censor del periódico El Siglo, que por eso no pudo publicar desde el 11 de julio hasta el 6 de agosto de 1944.\end{ejemplo}
\begin{cita}«La acción armada ejercida para capturar a Efraín González en cumplimiento de una orden judicial expedida por funcionario competente no constituye falla del servicio y fue, por lo mismo, legítima, pero ella causó un perjuicio económico a un tercero ajeno a esos hechos, consistente en la destrucción de una casa de propiedad de ese tercero, razón por la cual al Estado corresponde indemnizar el perjuicio causado, lo que equivale a hacer una equitativa distribución de las cargas públicas entre todos los contribuyentes desde luego que tal indemnización deberá hacerse con cargo al presupuesto de la Nación.»\end{cita}
\begin{ejemplo}Caso Efraín González (Consejo de Estado, 23 de mayo de 1973, Exp. 978): destrucción de una casa en la captura de Efraín González. El Ejército destruyó una casa para capturar a Efraín González, y la propietaria demandó al Estado. El demandante había arrendado la casa y quien la tenía en arriendo refugió a Efraín González. La acción armada, ejecutada en cumplimiento de una orden judicial de funcionario competente, fue legítima y no constituyó falla del servicio, pero causó un perjuicio a un tercero ajeno a los hechos: el Estado debía indemnizar con cargo al presupuesto de la Nación, como distribución equitativa de las cargas públicas entre todos los contribuyentes.\end{ejemplo}
\begin{cita}«Se puso en evidencia que la obra del puente de la 53 con la carrera 30 produjo un daño de carácter excepcional a los dueños del inmueble aledaño a dicha obra (número 28A-05 de la calle 53). Daño o perjuicio que no surge de una falla del servicio (la actividad de la entidad demandada fue legítima) sino del hecho de habérsele impuesto a los demandantes una carga especial en beneficio de la comunidad. Carga que rompe el principio de la igualdad frente a las cargas públicas (forma del principio general de la igualdad ante la ley).»\end{cita}
\begin{ejemplo}Caso del puente de la calle 53 con carrera 30 (Consejo de Estado, 30 de enero de 1987): la obra causó un daño excepcional a los dueños del inmueble aledaño. El daño no surgió de una falla del servicio, pues la actividad fue legítima, sino de una carga especial impuesta en beneficio de la comunidad, que rompió la igualdad frente a las cargas públicas.\end{ejemplo}
\section{La Asamblea Nacional Constituyente de 1991}
\subsection{Antecedentes}
La Constituyente de 1991 constitucionalizó el principio de la responsabilidad extracontractual del Estado. El artículo 90 se convirtió en la cláusula general de responsabilidad del Estado.

\subsection{La ponencia del artículo 90}
\textbf{1. De la antijuridicidad de la conducta a la antijuridicidad del daño.}

\begin{cita}«(…) En otras palabras, se desplaza el soporte de la responsabilidad administrativa, del concepto subjetivo de la antijuridicidad de la acción del Estado al concepto objetivo de la antijuridicidad del daño producido por ella. Esta antijuridicidad habrá de predicarse cuando se cause un detrimento patrimonial que carezca de título jurídico válido y que exceda el conjunto de las cargas que normalmente debe soportar el individuo en su vida social. (…)» (Gaceta Constitucional No. 56, p. 14)\end{cita}
Antes, especialmente bajo el régimen de falla del servicio, el análisis se concentraba en la conducta de la Administración: había que determinar si el Estado había actuado irregularmente, tardíamente, de manera deficiente o había omitido actuar. Con la Constitución de 1991, el centro del análisis se desplaza hacia la víctima y el daño sufrido: lo determinante es si existe un daño antijurídico, entendido como aquel que la persona no tiene el deber jurídico de soportar. Esto no significa que desaparezca la falla del servicio: el artículo 90 no adoptó un único régimen de responsabilidad; siguen existiendo títulos subjetivos, como la falla del servicio, y objetivos, como el daño especial y el riesgo excepcional.

\begin{comentario}Idea clave: antes el énfasis estaba en preguntar «¿el Estado actuó mal?»; a partir del artículo 90, el punto de partida es «¿la víctima sufrió un daño que jurídicamente no estaba obligada a soportar, y ese daño es imputable al Estado?».\end{comentario}
\textbf{2. Ampliación de la responsabilidad del Estado: responsabilidad por la actividad judicial y por el hecho del legislador.}

\begin{cita}«(…) Tal como se ha redactado el artículo, cabe perfectamente la posibilidad, hacia la cual claramente se está inclinando el derecho moderno, de extender el régimen de la responsabilidad patrimonial del Estado a aquella que se deriva de los yerros de la administración de justicia y eventualmente, en el futuro, también a la responsabilidad que pueda derivarse de la función legislativa. (…)» (Gaceta Constitucional No. 56, p. 14)\end{cita}
La cláusula constitucional permite comprender la responsabilidad estatal de manera general y no únicamente frente a la actividad administrativa: puede surgir por actuaciones u omisiones de las distintas autoridades públicas, incluidos supuestos relacionados con la actividad legislativa y la administración de justicia. La Corte Constitucional ha precisado además que el artículo 90 comprende ámbitos contractuales y extracontractuales (ver la objeción en la sección 3.4.3).

\textbf{3. Responsabilidad patrimonial de la entidad y del agente estatal.}

\begin{cita}«(…) Finalmente, he estimado de suma utilidad acoger dos planteamientos contenidos en algunos proyectos: el de que, para proteger a las víctimas, se les permita demandar indistintamente al Estado o a la autoridad pública involucrada; y el de que, para proteger el patrimonio público y combatir la indolencia personal del grueso de los funcionarios públicos, se le imponga al Estado el deber de repetir contra el funcionario culpable por las sumas por las que aquél hubiera sido condenado. (…)» (Gaceta Constitucional No. 56, p. 14)\end{cita}
Se acogieron dos propuestas: permitir a las víctimas demandar indistintamente al Estado o a la autoridad involucrada, y obligar al Estado a repetir contra el funcionario culpable, para proteger el patrimonio público y combatir la indolencia de los funcionarios. El modelo constitucional distingue ambas responsabilidades: frente a la víctima responde patrimonialmente el Estado cuando se cumplen los presupuestos del artículo 90; si el Estado es condenado y el daño fue consecuencia de la conducta dolosa o gravemente culposa de uno de sus agentes, surge el deber de ejercer la acción de repetición contra este.

\textbf{4. La culpa grave o el dolo como limitación para repetir.}

\begin{cita}«(…) El delegatario Ramírez Ocampo (…) insiste en reducir el artículo, agrega que la repetición contra el funcionario va a conducir a que no haya quien acepte desempeñar cargos públicos, el constituyente Uribe Vargas se muestra de acuerdo. (…)» (Gaceta Constitucional No. 132, p. 19)\end{cita}
Por eso la repetición quedó limitada al dolo o la culpa grave.

\subsection{El artículo 90 de la Constitución Política}
\begin{cita}«ARTÍCULO 90. El Estado responderá patrimonialmente por los daños antijurídicos que le sean imputables, causados por la acción o la omisión de las autoridades públicas. En el evento de ser condenado el Estado a la reparación patrimonial de uno de tales daños, que haya sido consecuencia de la conducta dolosa o gravemente culposa de un agente suyo, aquél deberá repetir contra éste.»\end{cita}
\textbf{Estructura.} El artículo 90 tiene dos incisos:

\begin{itemize}
\item \textbf{Primer inciso:} responsabilidad patrimonial del Estado frente a la víctima. De él salen los dos grandes elementos constitucionales de la responsabilidad del Estado: (1) el daño antijurídico y (2) la imputación al Estado.
\item \textbf{Segundo inciso:} responsabilidad del agente estatal frente al Estado mediante la repetición, cuando actuó con dolo o culpa grave.
\end{itemize}
\formula{Víctima → demanda al Estado → el Estado indemniza → si hubo dolo o culpa grave del agente → el Estado repite contra el agente}
\textbf{Alcance.} La Corte Constitucional tiene la posición de que el artículo 90 es una cláusula general que se aplica a toda la responsabilidad, extracontractual y contractual. Puede objetarse que no debería aplicarse a la responsabilidad contractual, porque allí no se miran daños antijurídicos sino el incumplimiento del contrato.

\begin{comentario}El primer inciso regula la responsabilidad patrimonial del Estado y el segundo la responsabilidad patrimonial del agente, es decir, la acción de repetición (ver Capítulo 10).\end{comentario}
\subsection{El daño antijurídico}
\begin{definicion}\textbf{Daño antijurídico:} el daño que la víctima no está en el deber jurídico de soportar; el que rompe el equilibrio ante las cargas públicas. No significa necesariamente que la actuación estatal sea ilegal.\end{definicion}
Para que el daño sea antijurídico se requieren tres elementos, cada uno frente a su opuesto:

{\small\begin{xltabular}{\linewidth}{|X|X|}\hline
\rowcolor{navy!12}\textbf{El daño antijurídico es…} & \textbf{…y se opone al daño que se debe soportar, que es…} \\\hline\endhead
Particular: afecta a una persona en concreto & General: una carga de todos, como los impuestos o el pico y placa \\\hline
Anormal o grave & Normal: por ejemplo, una revisión policial ordinaria (deja de serlo si, durante la revisión, el agente le rompe una pierna a la persona) \\\hline
Sin título jurídico que obligue a la víctima a soportarlo (inexistencia de título jurídico) & Con título jurídico válido: por ejemplo, una sentencia condenatoria válida justifica la privación de la libertad; no hay título si la persona fue privada de la libertad sin justa causa \\\hline
\end{xltabular}}
\begin{ejemplo}Caso del profesor privado de la libertad. Una menor fue víctima de un acto sexual en el colegio. Medicina Legal concluyó que hubo acto sexual y había una inferencia razonable en contra del profesor que justificaba detenerlo. El profesor duró siete meses privado de la libertad y luego fue absuelto por duda. ¿Sufrió un daño antijurídico?\end{ejemplo}
\textbf{Las dos posiciones en la Sección Tercera del Consejo de Estado:}

\begin{itemize}
\item \textbf{Subsecciones A y C:} para obtener indemnización debe acreditarse una falla del servicio, que se presenta, por ejemplo, cuando la medida de aseguramiento fue ilegal, irrazonable, desproporcionada o no cumplía los requisitos legales. Si la medida fue legal, en principio el Estado no responde bajo este título.
\item \textbf{Subsección B:} en algunos casos ha sostenido una posición de responsabilidad objetiva, aplicando el daño especial: aunque la medida de aseguramiento haya sido legal, puede haber responsabilidad si la persona finalmente absuelta soportó una carga particular y anormal que no estaba jurídicamente obligada a soportar.
\end{itemize}
\textbf{¿Por qué daño especial?} Porque, aunque el Estado pudo haber actuado lícitamente al imponer la medida de aseguramiento —en ese momento había elementos que permitían inferir razonablemente la participación del profesor—, finalmente fue absuelto y no se desvirtuó su presunción de inocencia; terminó soportando una carga particular, grave y anormal, superior a la que normalmente debe soportar cualquier ciudadano, lo que rompe el equilibrio frente a las cargas públicas, aun cuando no necesariamente haya existido una falla del servicio.

\begin{comentario}Sobre la evolución de la jurisprudencia en privación injusta de la libertad (unificaciones de 2006 y 2018, SU-072 de 2018) y el caso de la trata de personas, ver Capítulo 7, sección 7.3.\end{comentario}
\subsection{Responsabilidad del Estado y culpa personal del agente}
\textbf{Regla general.} Según el artículo 90, la víctima demanda al Estado y, si este resulta condenado porque su agente actuó con dolo o culpa grave, posteriormente ejerce la acción de repetición contra el agente. Los sujetos frente a la víctima de un daño antijurídico son tres: la víctima, la entidad estatal y el agente estatal.

\begin{definicion}\textbf{Culpa o falta personal del agente:} se presenta cuando el servidor público causa el daño actuando dentro de su órbita estrictamente personal o privada, sin ningún nexo con el servicio. En esos casos el daño no es imputable al Estado, porque la simple calidad de servidor público del autor no basta para comprometer la responsabilidad de la entidad.\end{definicion}
\begin{cita}El Consejo de Estado ha dicho que «la simple calidad de funcionario público que ostente el autor del hecho no vincula necesariamente al Estado», porque el funcionario puede actuar en su ámbito privado, completamente separado de la actividad pública.\end{cita}
\begin{ejemplo}Un soldado, estando fuera del servicio, se emborracha y, por una discusión personal relacionada con una mujer, usa un arma y mata a otra persona. Si el hecho obedeció exclusivamente a una venganza, riña o motivo personal y no tuvo relación con sus funciones, se configura una culpa personal del agente sin conexión con el servicio: el Estado no responde y la responsabilidad recae personalmente sobre el agente. Frente al Estado, ese hecho se trata como un hecho de un tercero que lo exonera de responsabilidad (ver Capítulo 5, sección 5.9).\end{ejemplo}
\textbf{Pregunta clave:} ¿el agente produjo el daño con ocasión o por razón del servicio, o actuó exclusivamente como un particular?

\begin{itemize}
\item \textbf{Si existe nexo con el servicio:} responde el Estado frente a la víctima y, si hubo dolo o culpa grave, puede repetir contra el agente.
\item \textbf{Si existe una culpa personal completamente desligada del servicio:} el daño no se imputa al Estado y la responsabilidad corresponde personalmente al agente; la demanda se dirige ante la jurisdicción ordinaria.
\end{itemize}
\subsection{Solidaridad y responsabilidad}
\begin{itemize}
\item \textbf{Responsabilidad:} se paga porque hay un daño y un nexo causal (una relación jurídica de imputación o causalidad).
\item \textbf{Solidaridad:} el pago se da sin relación de causalidad; es una ayuda económica.
\end{itemize}
\begin{comentario}No debe confundirse con las obligaciones solidarias (art. 2344 del C.C.), en las que varios deudores pueden estar llamados a responder por la totalidad frente al acreedor (ver Capítulo 5, secciones 5.2 y 5.11.1).\end{comentario}
\section{La responsabilidad extracontractual del Estado después de 1991}
¿Qué régimen de responsabilidad consagra el artículo 90? Se distinguen dos etapas:

{\small\begin{xltabular}{\linewidth}{|X|X|}\hline
\rowcolor{navy!12}\textbf{Período} & \textbf{Postura} \\\hline\endhead
1991-1993 & Plena aplicación del artículo 90 \\\hline
1993 a hoy & Regreso a los títulos de imputación preconstitucionales \\\hline
\end{xltabular}}
\subsection{Plena aplicación del artículo 90 (1991-1993)}
Inicialmente se consideró que el artículo 90 establecía un régimen objetivo de responsabilidad, porque desplazó el análisis desde la conducta del Estado hacia el daño antijurídico.

\begin{cita}«(…) es necesario dar aplicación al artículo 90 de la misma Carta que consagra la responsabilidad patrimonial del estado y la condiciona a la ocurrencia de un daño antijurídico que le sea imputable y que sea causado por la acción u omisión de las autoridades públicas. (…) Probados los dos elementos que exige la Constitución como necesarios para el compromiso de la responsabilidad estatal de carácter patrimonial, es irrelevante para ese solo efecto, que la acción dañosa sea constituida o no de falla del servicio; esa adquirirá importancia para el arreglo final de las cuentas entre la administración y el funcionario comprometido cuando esa circunstancia se presenta, que no es del caso. (…)» (Consejo de Estado, Sección Tercera, Expediente 5946, sentencia del 29 de octubre de 1992)\end{cita}
En síntesis: probados el daño antijurídico imputable y su causación por la acción u omisión de las autoridades, es irrelevante que la conducta constituya o no falla del servicio; la falla solo importa para el arreglo de cuentas entre la administración y el funcionario.

\begin{cita}«(…) se impone concluir que el centro de imputación jurídica demandado es responsable de los perjuicios causados con el incendio en la casa vecina con apoyo en el artículo 90 de la Constitución Nacional que en forma expresa dispone que el Estado debe responder "... patrimonialmente por los daños antijurídicos que le sean imputables, causados por la acción o la omisión de las autoridades públicas". Es esta responsabilidad, como reiteradamente lo ha dicho la Sala, la atención del constituyente no se fijó en el autor de la conducta causante del daño, sino en la víctima misma. Por ello importa más reparar el daño causado, que castigar una acción u omisión administrativa culpable. La finalidad de la responsabilidad patrimonial no consiste, pues, en BORRAR UNA CULPA, sino en hacer recaer sobre el patrimonio de la administración, el DAÑO SUFRIDO por el particular. La Culpa ha dejado de ser el FUNDAMENTO ÚNICO del sistema indemnizatorio, convirtiéndose simplemente en uno de los criterios jurídicos de imputación de daños a la administración. Esta verdad jurídica explica la razón por la cual en muchos casos ella deba responder de los daños patrimoniales ocasionados no obstante haber tenido una actuación lícita. La filosofía que informa todo este universo jurídico se apoya en el PRINCIPIO DE SOLIDARIDAD, que se recoge también en el artículo 1o. de la Constitución Nacional cuando se refiere a Colombia como un Estado Social de Derecho, fundado en el respeto de la dignidad de la persona humana y en LA SOLIDARIDAD DE LAS PERSONAS QUE LA INTEGRAN. (…)» (Consejo de Estado, Sección Tercera, Expediente 7130, sentencia del 26 de noviembre de 1992)\end{cita}
En síntesis (caso del incendio de la casa vecina): el constituyente fijó su atención en la víctima, no en el autor de la conducta; importa más reparar el daño que castigar una acción culpable. La culpa dejó de ser el fundamento único y pasó a ser uno de los criterios de imputación; por eso el Estado responde incluso por actuaciones lícitas. Su base es el principio de solidaridad del artículo 1.º de la Constitución.

\begin{cita}«(…) Se consagró entonces en la norma superior transcrita la responsabilidad patrimonial del Estado, por las acciones u omisiones de las autoridades públicas que ocasionen un daño antijurídico, régimen bajo el cual, para casos como el examinado, a la parte actora le corresponde demostrar que por parte de la autoridad respectiva se dio un comportamiento activo u omisivo, que le generó algún daño y que esa conducta y el daño respectivo se encuentran causalmente relacionados. Frente a esta situación el Estado podrá exonerarse de responsabilidad, solamente si acredita la fuerza mayor, el hecho exclusivo de la víctima o el hecho exclusivo de un tercero, sin que pueda alegar como eximentes el caso fortuito o una conducta diligente y cuidadosa, por cuanto en ese tipo de responsabilidad, antes que la antijuridicidad del actuar del ente administrativo, lo que importa es la antijuridicidad del daño, el cual, resulta antijurídico cuando a pesar de ser legítima la conducta de la administración, la misma causa un perjuicio a quien no estaba obligado a soportarlo. (…)» (Consejo de Estado, Sección Tercera, Expediente 7340, sentencia del 17 de febrero de 1993)\end{cita}
En síntesis: la parte actora debe demostrar la conducta activa u omisiva de la autoridad, el daño y la relación causal. El Estado solo se exonera con la fuerza mayor, el hecho exclusivo de la víctima o el hecho exclusivo de un tercero; no puede alegar el caso fortuito ni su diligencia, porque importa la antijuridicidad del daño y no la de la conducta.

\subsection{Regreso a los títulos de imputación preconstitucionales (1993 a hoy)}
\begin{cita}«(…) No es del todo exacto que luego de la expedición de la Carta de 1991 la responsabilidad estatal se volvió objetiva y que en ningún evento se pueda probar la conducta irregular de la administración que produjo el daño, aunque sí puede estimarse que la jurisprudencia deberá tener también en cuenta ese criterio objetivista para su interpretación. (…) En síntesis, la nueva Constitución, a pesar de su amplitud en materia de responsabilidad, no la hizo exclusivamente objetiva ni borró del ordenamiento la responsabilidad por falla del servicio, ni por cualquiera otra fuente de las aceptadas por el derecho administrativo. Las nociones de imputabilidad y de daño antijurídico así lo dan a entender. (…)» (Consejo de Estado, Sección Tercera, Expediente 7429, sentencia del 2 de marzo de 1993)\end{cita}
A partir de 1993, el Consejo de Estado reinterpretó la norma: el artículo 90 no volvió objetiva toda la responsabilidad del Estado ni eliminó la falla del servicio ni las demás fuentes aceptadas por el derecho administrativo; las nociones de daño antijurídico e imputación permiten mantener los distintos títulos de imputación.

\begin{itemize}
\item Por la vía de la imputación, el Consejo de Estado volvió a aplicar los tres títulos tradicionales: falla del servicio, riesgo excepcional y daño especial.
\item La falla del servicio es el principal título de imputación y corresponde a un régimen subjetivo, porque exige analizar la conducta irregular de la Administración. Excepcionalmente se aplican el riesgo excepcional y el daño especial, que corresponden a regímenes objetivos, en los que no es necesario demostrar que el Estado actuó mal.
\end{itemize}
\begin{comentario}Crítica: esta reinterpretación puede cuestionarse, porque ni el texto del artículo 90 ni los debates de la Gaceta Constitucional dicen expresamente que deban mantenerse esos tres regímenes. El artículo 90 simplemente establece que el Estado responde por los daños antijurídicos que le sean imputables; fue la jurisprudencia la que, mediante el concepto de imputación, reincorporó la falla del servicio, el riesgo excepcional y el daño especial.\end{comentario}
\textbf{Síntesis: la responsabilidad del Estado en la Constitución de 1991} El artículo 90 establece una cláusula general de responsabilidad patrimonial que se estructura sobre dos elementos:

\begin{itemize}
\item \textbf{Daño antijurídico:} una lesión que la víctima no tiene el deber jurídico de soportar; no significa necesariamente que la actuación estatal sea ilegal.
\item \textbf{Imputación al Estado:} debe demostrarse que el daño es atribuible a una acción u omisión de una autoridad pública, bajo alguno de los títulos:
\begin{itemize}
\item Falla del servicio: régimen subjetivo; se acredita el incumplimiento de un deber estatal.
\item Riesgo excepcional: régimen objetivo; el daño proviene de la materialización de un riesgo excepcional atribuible al Estado.
\item Daño especial: régimen objetivo; una actuación estatal lícita impone a la víctima una carga especial y anormal, rompiendo la igualdad ante las cargas públicas.
\end{itemize}
\end{itemize}
La Constitución de 1991 no convirtió toda la responsabilidad estatal en objetiva ni eliminó la falla del servicio: estableció como base el daño antijurídico y su imputación, bajo distintos títulos.

\section{\emph{Iura novit curia} y estudio de los títulos de imputación}
Es necesario distinguir entre el principio de congruencia, el principio de jurisdicción rogada y el \emph{iura novit curia}.

\subsection{El principio de congruencia}
\begin{definicion}\textbf{Principio de congruencia:} exige correspondencia entre lo pedido por las partes y lo decidido por el juez.\end{definicion}
\textbf{Clases de incongruencia:}

\begin{itemize}
\item \textbf{Extra petita:} el juez concede algo distinto de lo solicitado.
\end{itemize}
\begin{ejemplo}Se pidieron únicamente daño emergente y lucro cesante (o solo daño emergente), pero la sentencia reconoce perjuicios morales no solicitados.\end{ejemplo}
\begin{itemize}
\item \textbf{Ultra petita:} el juez concede más de lo pedido.
\end{itemize}
\begin{ejemplo}El demandante solicita \$100 millones por perjuicios y el juez reconoce \$150 millones.\end{ejemplo}
\begin{itemize}
\item \textbf{Citra o infra petita:} el juez deja sin resolver lo que debía decidir.
\end{itemize}
La congruencia aplica con mayor rigor en los procesos donde se discuten derechos dispositivos. Se flexibiliza en algunos procesos, como los laborales, la tutela y las acciones de grupo (ciertos procesos colectivos), donde el juez puede tener facultades más amplias de protección según la naturaleza de los derechos involucrados.

\subsection{El principio de jurisdicción rogada}
\begin{definicion}\textbf{Principio de jurisdicción rogada:} el juez, especialmente en lo contencioso administrativo, solo puede pronunciarse sobre aquello que las partes le solicitan y dentro de los cargos planteados en la demanda.\end{definicion}
Aplica especialmente en los medios de control contra actos administrativos, en los que el juez analiza la legalidad del acto con base en los cargos de violación formulados por el demandante.

\begin{ejemplo}Sale una reglamentación en materia de contratación estatal y se demanda su artículo 2. Por el principio de jurisdicción rogada, el juez, en principio, solo puede anular el artículo 2 (no el 5) y con los motivos alegados.\end{ejemplo}
\textbf{Excepciones (no aplica o se flexibiliza):}

\begin{itemize}
\item \textbf{Nulidad por inconstitucionalidad (art. 135 del CPACA):} el juez puede anular por cargos distintos y normas que no fueron demandadas si considera que forman una unidad.
\item \textbf{Control inmediato de legalidad:} el control puede ser más amplio u oficioso.
\end{itemize}
\subsection{El principio \emph{iura novit curia}}
\begin{definicion}\textbf{Iura novit curia:} «el juez conoce el derecho». «Incumbe a las partes dar los hechos y al juez el derecho aplicable». El juez puede aplicar la norma o el título jurídico correcto aunque las partes hayan citado otro o ninguno, siempre que respete los hechos de la demanda, las pretensiones y el derecho de defensa.\end{definicion}
\textbf{Reglas:}

\begin{itemize}
\item El juez no está atado al título de imputación invocado en la demanda.
\item No puede modificar los hechos (la imputación fáctica) ni las pretensiones.
\item Puede hacer un estudio sucesivo de los títulos de imputación.
\end{itemize}
\begin{ejemplo}En una reparación directa, el demandante sostiene que hubo falla del servicio, pero el juez concluye que, con los mismos hechos probados, el título correcto es el riesgo excepcional. Puede cambiar el título jurídico sin violar el principio de congruencia, porque no modifica los hechos, sino únicamente su calificación jurídica.\end{ejemplo}
\begin{ejemplo}Caso de Pedro y Juan. Pedro y Juan eran ganaderos y no regresaron a su casa; al día siguiente aparecieron muertos. Sus familiares demandaron al Estado afirmando que habían sido víctimas de un falso positivo: que fueron presentados como muertos en combate dentro de un operativo que, según los demandantes, no había existido. Durante el proceso se probó que sí había existido un operativo militar, de modo que no quedó demostrada la imputación fáctica planteada en la demanda. Si el juez, pese a ello, cambia los hechos y condena al Estado partiendo de una versión distinta —por ejemplo, que murieron durante un operativo real, pero por un uso desproporcionado de la fuerza—, ya no estaría aplicando el \emph{iura novit curia}: no estaría cambiando solo el derecho aplicable, sino también la causa fáctica de la demanda.\end{ejemplo}
\begin{ejemplo}Caso de los niños del accidente del Agustiniano. Sirve para estudiar cómo el Consejo de Estado puede escoger el título de imputación aplicable sin alterar la base fáctica. Al condenar hubo una dualidad de título de imputación y se condenó por falla del servicio «por la función».\end{ejemplo}
\begin{comentario}Idea clave: el \emph{iura novit curia} permite al juez cambiar la calificación jurídica, pero no inventar o sustituir los hechos que fundamentan la demanda. El juez es libre de escoger el título de imputación, pero no los argumentos fácticos.\end{comentario}
\subsection{Diferencias entre los tres principios}
{\small\begin{xltabular}{\linewidth}{|X|X|}\hline
\rowcolor{navy!12}\textbf{Principio} & \textbf{Qué controla} \\\hline\endhead
Jurisdicción rogada & Qué puede estudiar el juez (lo pedido y los cargos planteados) \\\hline
Congruencia & La correspondencia entre lo pedido, lo debatido y lo decidido \\\hline
Iura novit curia & La libertad del juez para escoger el derecho aplicable, sin cambiar los hechos \\\hline
\end{xltabular}}
\subsection{Oposición y excepciones}
\begin{definicion}\textbf{Oposición:} es la defensa general del demandado frente a la demanda. Puede consistir en negar los hechos, controvertir las pruebas, discutir la interpretación jurídica o rechazar las pretensiones.\end{definicion}
\textbf{Tipos de oposición:}

\begin{itemize}
\item \textbf{Fáctica:} se niegan o discuten los hechos.
\item \textbf{Probatoria:} se controvierten las pruebas.
\item \textbf{Jurídica:} se discute la consecuencia legal que pretende el demandante.
\end{itemize}
\begin{definicion}\textbf{Excepciones:} son defensas jurídicas concretas, dentro de la oposición, dirigidas a impedir, modificar o extinguir la pretensión del demandante.\end{definicion}
\textbf{Tipos de excepciones:}

\begin{itemize}
\item \textbf{Previas:} cuestionan aspectos procesales, como la competencia, la indebida representación o la ineptitud de la demanda.
\item \textbf{De mérito o de fondo:} atacan directamente el derecho reclamado; proponen un hecho nuevo que busca destruir la demanda. Ejemplos: prescripción, pago, inexistencia de la obligación, culpa exclusiva de la víctima, hecho de un tercero o fuerza mayor.
\end{itemize}
\begin{comentario}En resumen: la oposición es el género; la excepción es una especie de defensa dentro de la oposición.\end{comentario}
\begin{comentario}En materia de error judicial, el juez tampoco puede cambiar la imputación fáctica: si el demandante afirma que la fuente del daño es una providencia judicial, el caso debe estudiarse como error judicial; el juez no puede modificar los hechos ni «mover de título» el caso a partir de hechos distintos. Sobre el ejemplo de los antibióticos y el conocimiento privado del juez, ver Capítulo 5 (responsabilidad médica).\end{comentario}
\section{Función de la responsabilidad extracontractual del Estado}
{\small\begin{xltabular}{\linewidth}{|X|X|}\hline
\rowcolor{navy!12}\textbf{Fundamento} & \textbf{Función} \\\hline\endhead
Falla del servicio & Función demarcatoria \\\hline
Daño antijurídico & Función indemnizatoria \\\hline
\end{xltabular}}
\section{Caso integrador: vehículo oficial y actividad peligrosa (el alcalde de Pinchote)}
\begin{ejemplo}Hechos. El alcalde del municipio de Pinchote conduce un vehículo de propiedad del municipio, en desarrollo de una actividad misional. Sufre un infarto, pierde el control del vehículo y ocasiona un accidente en el que mueren tres peatones.\end{ejemplo}
\begin{itemize}
\item \textbf{Título de imputación:} en principio puede existir responsabilidad patrimonial del municipio bajo el título de riesgo excepcional, porque la conducción de vehículos es una actividad peligrosa, el vehículo pertenece al municipio y existe nexo con una actividad misional.
\item \textbf{Exoneración:} el infarto no exonera automáticamente; debe analizarse si constituye una causa extraña con aptitud para romper la imputación.
\item \textbf{Acción:} medio de control de reparación directa (art. 140 del CPACA). Los familiares o personas legitimadas demandan al municipio de Pinchote y solicitan la reparación de los perjuicios derivados del daño antijurídico.
\end{itemize}
\begin{comentario}Ejemplo análogo a propósito del requisito de exterioridad de la causa extraña: si mientras conduzco me da un paro cardíaco y atropello a un peatón, no me exonero, porque la salud del conductor no es externa al ámbito del deudor; en últimas, se responde por el riesgo (ver Capítulo 5, sección 5.8.3).\end{comentario}
